\section{An explicit three-adic kernel class}
\label{sec:explicit-core}
\subsection{The ring and the class}\label{c3:sec:introduction}

\medskip
\noindent\textbf{Generic restriction.}
Products of units give explicit classes with which to study restriction
from a regular local ring to its fraction field. Dennis--Stein identities
can establish the vanishing of such a class after localization
\cite[III, Definition~5.11 and Theorem~5.11.1]{Weibel}.
For nonvanishing, we use the continuous fiber comparison
\cite[Proposition~4.10]{AMMN}, together with multiplicative traces and
the module structure on cyclic homology
\cite{BGTTrace,NikolausScholze,HoyoisCircle,BokstedtOttosen}.
The example below requires both an integral calculation in algebraic
$\CthreeK$-theory and an ordinary, possibly non-Hausdorff, differential quotient.

\medskip
\noindent\textbf{An explicit class.}
For units $b_1,\ldots,b_r$ in a commutative ring $R$, write
$\{b_1,\ldots,b_r\}=[b_1]\cdots[b_r]\in\CthreeK_r(R)$ for the product of their
classes in actual Quillen algebraic $\CthreeK$-theory.

\begin{theorem}
\label{thm:explicit-three-dimensional}
Let
\[
 A_3=
 \left(\CthreeZ_{(3)}[x,y,z]/(3+xy+z^2)\right)_{(x,y,z)},
 \qquad
 \widehat A_3=\CthreeZ_3[[x,y,z]]/(3+xy+z^2).
\]
Then $\widehat A_3$ is the maximal-ideal completion of $A_3$, and both
rings are ramified regular local domains of dimension three.
For either $R=A_3$ or $R=\widehat A_3$, the class
\[
 c_R=
 2\left\{\frac{1+z}{2},\,1+x,\,
              1+\frac{(1+x)y}{4}\right\}\in\CthreeK_3(R)
\]
has infinite order and maps to zero in $\CthreeK_3(R[1/x])$ integrally.
Consequently,
\[
 \CthreeK_3(R)\otimes_{\CthreeZ}\CthreeQ
 \longrightarrow
 \CthreeK_3(\CthreeFrac(R))\otimes_{\CthreeZ}\CthreeQ
\]
is not injective.
\end{theorem}

The same unit formula therefore gives a nonzero rational class in both
the algebraic local ring and its completion. Its generic vanishing
already holds before rationalization.

\medskip
\noindent\textbf{Method.}
The integral calculation uses a Dennis--Stein symbol whose product with
twice a unit class is $c_R$.
For nonvanishing, we map to a smooth formal quadric over an Eisenstein
quadratic extension of $\CthreeZ_3$.
Proposition~\ref{c3:prop:relative-character} constructs the relative
character with its $\CthreeK$-module action and fixes its value on relative
units. Lemma~\ref{c3:lem:ordinary-residue} then detects its top component by
two endpoint values, without constructing a convergent primitive for a
closed one-form. The endpoint difference is nonzero by an estimate
valid for every term of its defining series.
Vanishing of the actual fourth $\CthreeK$-group of the special fiber finally
passes this detection from relative to absolute $\CthreeK$-theory.

Section~\ref{c3:sec:preliminaries} fixes the conventions and the continuous
comparison. Section~\ref{c3:sec:integral} proves integral vanishing.
Section~\ref{c3:sec:character} constructs and normalizes the character.
Section~\ref{c3:sec:residue} proves the residue criterion.
Section~\ref{c3:sec:nonvanishing} completes the proof.







\subsection{Preliminaries}\label{c3:sec:preliminaries}

We use Quillen algebraic $\CthreeK$-theory and its unit products with the
conventions of \cite{Quillen,Weibel}.
For continuous cyclic constructions we follow
\cite[Definition~4.1 and Proposition~4.2]{AMMN}.
We write $\mathrm{HH}$, $\mathrm{HC}$, $\mathrm{HC}^{-}$, and $\mathrm{HP}$ for Hochschild, cyclic, negative cyclic, and periodic cyclic homology, respectively. The symbols $\mathrm{THH}$ and $\mathrm{TC}$ denote topological Hochschild and topological cyclic homology.

\begin{notation}\label{c3:notation:k-theory}
An uncompleted group $\CthreeK_i(D)$ always means an actual Quillen
$\CthreeK$-group.
For an ideal $J\subseteq D$, the relative spectrum is
\[
 \CthreeK(D,J)=\Cthreefib\bigl(\CthreeK(D)\longrightarrow\CthreeK(D/J)\bigr),
 \qquad
 \CthreeK_i(D,J)=\pi_i\CthreeK(D,J).
\]
Thus relative $\CthreeK$-theory is defined by a homotopy fiber, rather than
by the kernel of a map on homotopy groups.

For a spectrum-valued invariant, the notation $;\CthreeQ_p$ means
$p$-completion of the spectrum followed by inversion of $p$.
It does not mean tensoring its actual homotopy groups with $\CthreeQ_p$.
The superscript $\Cthreects$ denotes the homotopy inverse limit over
the $p$-power reductions.
\end{notation}

Unit products are natural under ring maps, multilinear, and graded
commutative. They satisfy $\{b,1-b\}=0$ whenever both entries are units;
see \cite[III, Lemma~5.10.2, Definition~5.11, and the unit-product
discussion following Theorem~5.11.1]{Weibel}.
The distinction in Notation~\ref{c3:notation:k-theory} will be retained
when these products are mapped to continuous invariants.

\begin{setup}\label{c3:setup:coefficient-ring}
Put
\[
 W=\CthreeZ_3[u]/(u^2+3),\qquad E=W[1/3].
\]
The Eisenstein equation makes $W$ a complete discrete valuation ring
with uniformizer $u$ and residue field $\CthreeF_3$.
Its $u$-adic and $3$-adic topologies agree.
For a smooth $W$-algebra $D_0$ of relative dimension at most two, let
$D$ denote its $3$-adic completion, and put
\[
 \Omega_D^j=
 \left(\varprojlim_n\Omega^j_{D_0/W}/3^n\Omega^j_{D_0/W}\right)[1/3].
\]
All differentials on these modules are relative to $W$.
The notation $W\langle T_1,\ldots,T_r\rangle$ denotes the restricted
power series algebra, namely the $3$-adic completion of
$W[T_1,\ldots,T_r]$.

The two cases used below are $D=W\langle T\rangle$ and
\[
 B=W\langle X,Y,Z\rangle/(XY+Z^2-1).
\]
The algebra $W[X,Y,Z]/(XY+Z^2-1)$ is smooth of relative dimension two:
its partial derivatives $Y,X,2Z$ have no simultaneous zero on the
quadric.
\end{setup}

The continuous comparison will be used in its formal-scheme form.

\begin{proposition}[Continuous fiber comparison
{\cite[Proposition~4.10]{AMMN}}]\label{c3:prop:continuous-comparison}
Let $\CthreecO$ be the integer ring of a complete discretely valued
characteristic-zero field with perfect residue field $k$ of
characteristic $p$. Let $\mathfrak X$ be a smooth $p$-adic formal
$\CthreecO$-scheme. There is a natural diagram
\begin{equation}\label{c3:eq:continuous-comparison}
\begin{CD}
 \CthreeK^\Cthreects(\mathfrak X;\CthreeQ_p)
     @>>> \CthreeK(\mathfrak X_k;\CthreeQ_p)\\
 @VVV @VVV\\
 \CthreeTC^\Cthreects(\mathfrak X;\CthreeQ_p)
     @>>> \CthreeTC(\mathfrak X_k;\CthreeQ_p)\\
 @VVV @VVV\\
 \CthreeHC^{-,\Cthreects}(\mathfrak X/\CthreecO;\CthreeQ_p)
     @>>> \CthreeHP^\Cthreects(\mathfrak X/\CthreecO;\CthreeQ_p)
\end{CD}
\end{equation}
in which both squares are cartesian.
\end{proposition}

The proof of \cite[Proposition~4.10]{AMMN} identifies the horizontal
fibers after projection from absolute Hochschild homology to
Hochschild homology relative to $\CthreecO$.
For an algebraization with bounded $p$-power torsion,
\cite[Proposition~4.2]{AMMN} identifies continuous Hochschild and cyclic
constructions with their $p$-completed counterparts.
The relative version is stated immediately before
\cite[Proposition~4.10]{AMMN}.
These assertions apply to the smooth, $W$-flat algebraizations in
Setup~\ref{c3:setup:coefficient-ring}.

\subsection{Integral vanishing}\label{c3:sec:integral}

In this section, we prove the regularity and integral vanishing
assertions of Theorem~\ref{thm:explicit-three-dimensional}.

\begin{proposition}\label{c3:prop:integral}
The rings $A_3$ and $\widehat A_3$ of
Theorem~\ref{thm:explicit-three-dimensional} are ramified regular local
domains of dimension three, with $\widehat A_3$ the maximal-ideal
completion of $A_3$. For either ring $R$, all entries defining $c_R$
are units, and
\[
 c_R|_{R[1/x]}=0\quad\text{in }\CthreeK_3(R[1/x]).
\]
\end{proposition}

\begin{proof}
The ambient local ring
$\CthreeZ_{(3)}[x,y,z]_{(3,x,y,z)}$ is regular of dimension four.
The image of $3+xy+z^2$ in its maximal ideal modulo its square is the
nonzero image of $3$.
Its quotient is therefore regular of dimension three, with maximal
ideal $(x,y,z)$, and is a domain.
Completion gives the displayed presentation of $\widehat A_3$.
The same parameter calculation in $\CthreeZ_3[[x,y,z]]$ proves that this
completion is regular of dimension three.
In both rings, $3=-xy-z^2$ belongs to the square of the maximal ideal,
so the rings are ramified.
Each entry in the definition of $c_R$ has nonzero image in $\CthreeF_3$.

We use Dennis--Stein symbols with the convention that
$\langle r,s\rangle$ is defined when $1-rs$ is invertible.
In additive notation their relations are
\begin{align*}
 \langle r,s\rangle&=-\langle s,r\rangle,\\
 \langle r,s\rangle+\langle r,t\rangle
   &=\langle r,s+t-rst\rangle,\\
 \langle r,st\rangle
   &=\langle rs,t\rangle+\langle tr,s\rangle
\end{align*}
whenever the symbols involved are defined
\cite[III, Definition~5.11, relations (D1)--(D3)]{Weibel}.
If $r$ is a unit, then
$\langle r,s\rangle=\{r,1-rs\}$.
Skew-symmetry gives
$\langle r,s\rangle=\{1-rs,s\}$ when $s$ is a unit.
In particular, $\langle r,1\rangle=0$ whenever it is defined.
These identities hold in actual $\CthreeK$-groups.

In $R$, put
\[
 a=\frac{1+z}{2},\qquad t=\frac y4,\qquad
 \lambda=1+x,\qquad \mu=1+\lambda t,\qquad q=a(1-a).
\]
The defining equation of $R$ gives
\[
 q=\frac{1-z^2}{4}=1+xt.
\]
The elements $a$, $1-a$, $\lambda$, $\mu$, $q$, and $\lambda q$
are units. The second Dennis--Stein relation gives
\[
 \delta\coloneqq \langle -x,t\rangle
   =\langle -x,t\rangle+\langle -x,1\rangle
   =\langle -x,t+1+xt\rangle
   =\langle -x,\mu\rangle
   =\{\lambda q,\mu\},
\]
where the last equality uses $1+x\mu=\lambda q$.

Graded commutativity gives $2\{a,a\}=0$, and the Steinberg relation
gives $\{a,1-a\}=0$. Hence
\[
 2\{a,q\}=2\{a,a\}+2\{a,1-a\}=0.
\]
Multiplying $\delta$ by $2[a]$ now yields
\[
 2[a]\delta
   =2\{a,\lambda,\mu\}+2\{a,q,\mu\}
   =c_R.
\]
After inverting $x$, the first-entry formula for the Dennis--Stein
symbol identifies $\delta$ with $\{-x,q\}$. Consequently,
\[
 c_R|_{R[1/x]}
   =2\{a,-x,q\}
   =-2\{a,q,-x\}
   =0.
\]
No completion or rationalization of a $\CthreeK$-group enters this calculation.
\end{proof}

\subsection{The relative character and its normalization}
\label{c3:sec:character}

The goal of this section is to construct the character used to detect
$c_R$ and to determine its value on a relative unit multiplied by two
absolute unit classes.

\begin{proposition}
\label{c3:prop:relative-character}
In Setup~\ref{c3:setup:coefficient-ring}, the continuous comparison defines
a natural map from the actual relative spectrum
\[
 \chi_D\colon
 \CthreeK(D,uD)\longrightarrow
 \Sigma\CthreeHC^\Cthreects(D/W;\CthreeQ_3).
\]
This map is $\CthreeK(D)$-linear, where the target action is obtained by
restriction along
\[
 \CthreeK(D)\longrightarrow\CthreeHC^{-,\Cthreects}(D/W;\CthreeQ_3).
\]
For $v\in1+uD$, let $[v]_{\Cthreerel}\in\CthreeK_1(D,uD)$ be the relative unit
class specified by the scalar automorphism $v$ and its identity
reduction modulo $u$. Then
\[
 \chi_D([v]_{\Cthreerel})=\log v
 \quad\text{in }\CthreeHC^\Cthreects_0(D/W;\CthreeQ_3)=\Omega_D^0.
\]
For $b_1,b_2\in D^\times$, its product character satisfies
\begin{equation}\label{c3:eq:unit-character}
 \chi_D\bigl([v]_{\Cthreerel}[b_1][b_2]\bigr)_{\Cthreetopcomp}
   =\pm\log(v)\,\Cthreedlog b_1\wedge\Cthreedlog b_2
   \quad\text{in }\Omega_D^2/\Cthreedd\Omega_D^1.
\end{equation}
Here the top component is the first summand of
\begin{equation}\label{c3:eq:cyclic-degree-two}
 \CthreeHC^\Cthreects_2(D/W;\CthreeQ_3)
   =
   \bigl(\Omega_D^2/\Cthreedd\Omega_D^1\bigr)
   \oplus
   \ker\bigl(\Cthreedd\colon\Omega_D^0\longrightarrow\Omega_D^1\bigr).
\end{equation}
The sign in \eqref{c3:eq:unit-character} is the universal suspension sign.
Equation~\eqref{c3:eq:cyclic-degree-two} uses the ordinary differential image.
\end{proposition}

\begin{proof}
The character is obtained from the horizontal fibers of
\eqref{c3:eq:continuous-comparison} for $\CthreeSpf D$.
Compatible maps from actual to continuous $\CthreeK$-theory, followed by
completion of spectra and rationalization, give a map from
$\CthreeK(D,uD)$ to the fiber of the top row.
The two squares carry this fiber to the fiber of
$\CthreeHC^{-,\Cthreects}\to\CthreeHP^\Cthreects$.
For a circle object $M$, the norm sequence is
\[
 \Sigma M_{hS^1}\longrightarrow M^{hS^1}
       \longrightarrow M^{tS^1};
\]
see \cite[\S I.3 and Corollary~I.4.3]{NikolausScholze}.
By \cite[Theorem~2.1, Lemma~2.2, and Theorem~2.3]{HoyoisCircle},
its cyclic interpretation identifies the last fiber with
\[\Sigma\CthreeHC^\Cthreects(D/W;\CthreeQ_3).\]
This defines a map out of actual relative $\CthreeK$-theory.

The required module action can be followed through the construction of
the comparison. Put $H=\CthreeTHH(D)$, regarded as a commutative algebra in
cyclotomic spectra.
For a bounded-below cyclotomic spectrum $H$, the construction in
\cite[Proposition~2.8]{AMMN} tensors the cyclotomic maps
$\CthreeZ^{\Cthreetriv}\to\CthreeTHH(\CthreeF_p)$ with $H$.
It gives a cartesian square with $\CthreeTC$ in the top row and $\CthreeTC^-$ in
the bottom row.
Its horizontal fiber is the $p$-completion of
$\Sigma(H\otimes\CthreeZ_{hC_p})_{hS^1}$.
The comparison in \cite[Corollary~2.10]{AMMN} replaces the lower-right
term by a Tate construction, using the equivalence induced by
$\CthreeZ^{\Cthreetriv}\to\CthreeTHH(\CthreeF_p)$, and then inverts $p$.
Specializing to $H=\CthreeTHH(D)$ and using
$\CthreeTHH(D)\otimes\CthreeZ\to\CthreeHH(D)$ gives the square with rows
$\CthreeTC(D)\to\CthreeTC(D\otimes_{\CthreebS}\CthreeF_p)$ and
$\CthreeHC^-(D)\to\CthreeHP(D)$ in \cite[Theorem~2.12]{AMMN}.
The reduction construction of \cite[Definition~2.14]{AMMN} replaces
derived reduction by ordinary reduction through their rational
$\CthreeTC$-equivalence.
Finally, \cite[Proposition~4.10]{AMMN} replaces reduction modulo $p$
by reduction modulo a uniformizer and projects to relative
Hochschild homology.

Every object $H\otimes T$ in the first two constructions is an
$H$-module through its first factor, and every map
$\Cthreeid_H\otimes f$ is $H$-linear.
The lax monoidal structure of $\CthreeTC$ therefore makes their $\CthreeTC$ images
modules over $\CthreeTC(H)=\CthreeTC(D)$.
The homotopy-fixed-point and Tate objects carry actions of $\CthreeTC^-(H)$,
and hence of $\CthreeTC(H)$ through its canonical map to $\CthreeTC^-(H)$.
Their norm and comparison maps preserve these actions.
This compatibility follows either from the lax monoidal constructions
in \cite[\S I.3 and Corollary~I.4.3]{NikolausScholze}, or by forming the norm
inside the category of modules over the circle-equivariant algebra.

In particular, the Tate equivalence inverted in
\cite[Corollary~2.10]{AMMN} is a module map.
Its inverse is taken in the module category.
The map
$\CthreeTC(H)\to\CthreeTC(H\otimes\CthreeZ^{\Cthreetriv})$, which becomes an equivalence after
rationalization in \cite[Theorem~2.12]{AMMN}, is a module map induced
by the unit of $\CthreeZ^{\Cthreetriv}$.
After forgetting to circle objects, the maps
\[
 H\otimes\CthreeZ\longrightarrow\CthreeHH(D)
       \longrightarrow\CthreeHH(D/W)
\]
are maps of $H$-algebras.
The reduction maps
$D\otimes_{\CthreebS}\CthreeF_3\to D/3\to D/u$ are maps of $D$-algebras.
Thus the rational $\CthreeTC$-equivalences used in
\cite[Definition~2.14 and Proposition~4.10]{AMMN} are
$\CthreeTC(D)$-linear as well.

Completion, rationalization, and passage to fibers preserve these
module constructions. The resulting action on the cyclic fiber is
restriction along
\[
 \CthreeTC(D)\longrightarrow\CthreeTC^-(D)
       \longrightarrow\CthreeHC^{-,\Cthreects}(D/W;\CthreeQ_3).
\]
The multiplicative cyclotomic trace, whose composite with $\CthreeTHH$ is
the topological Dennis trace
\cite[Theorems~7.3 and~7.4]{BGTTrace}, and the compatible
actual-to-continuous maps make $\chi_D$ a $\CthreeK(D)$-module map.
The suspension contributes the usual sign for a suspended module.
Only the negative cyclic action on cyclic homology is used here;
no ring structure on $\CthreeHC$ is asserted.

For the calculation of the target, the relative dimension bound gives
a finite mixed-complex model. On the algebraization $D_0$, the
normalized Hochschild map is
\begin{equation}\label{c3:eq:finite-hkr}
 b_0\otimes\cdots\otimes b_n
   \longmapsto
   \frac{1}{n!}b_0\,\Cthreedd b_1\wedge\cdots\wedge\Cthreedd b_n
   \qquad(0\leq n\leq2),
\end{equation}
and is zero in higher degrees.
The factorials occurring here are units in $W$.
The Hochschild differential maps to zero, while the Connes operator
maps to $\Cthreedd$; beyond degree two its target is zero.
Smooth Hochschild--Kostant--Rosenberg identifies the underlying
Hochschild homology with $\Omega^n_{D_0/W}$, vanishing for $n>2$.
This identification follows from the Koszul resolution of the smooth
diagonal and is realized by \eqref{c3:eq:finite-hkr}.
Hence \eqref{c3:eq:finite-hkr} is a quasi-isomorphism of mixed complexes
to $(\Omega^\ast_{D_0/W},0,\Cthreedd)$.
Only this finite mixed complex is needed, rather than an infinite
rational splitting.

The cyclic, negative cyclic, and periodic total complexes of this
model contain only finitely many form modules in each total degree.
Those modules are $W$-flat, and their reduction towers modulo $3^n$
are termwise surjective.
Derived $3$-completion is therefore given by termwise completion
of these total complexes.
The relative continuity assertion of
\cite[Proposition~4.2 and the paragraph preceding
Proposition~4.10]{AMMN} identifies them with the continuous
constructions. Inverting $3$ gives
$\CthreeHC^\Cthreects_0(D/W;\CthreeQ_3)=\Omega_D^0$ and
\eqref{c3:eq:cyclic-degree-two}.
Indeed, cyclic degree two has $\Omega_D^2$ in column zero and
$\Omega_D^0$ in column minus one.
The differential out of the second term is $\Cthreedd$, and boundaries
into the first term form $\Cthreedd\Omega_D^1$.
These are homology groups in the ordinary derived category, with
no closure taken on a differential image.

A unit $v\in1+uD$ determines a particular relative class: its scalar
automorphism is a loop whose reduction is the identity automorphism.
This construction is natural in $D$ and maps to $[v]\in\CthreeK_1(D)$.
To determine its character, first take $D=W\langle T\rangle$ and
$v=1+uT$, and write
\[
 f=\chi_D([1+uT]_{\Cthreerel})\in E\langle T\rangle.
\]
Choose the circle orientation for which the norm is Connes' operator,
as in \cite[Lemma~2.2 and Theorem~2.3]{HoyoisCircle}.
The norm followed by projection to $\CthreeHH_1$ then sends $f$ to $\Cthreedd f$.
On the $\CthreeK$-theory side, this composite is the Dennis trace.
The scalar-unit calculation for that trace represents the class of
$b$ by $b^{-1}\otimes b$ in $\CthreeHH_1$, which
\eqref{c3:eq:finite-hkr} sends to $\Cthreedlog b$.
It follows that
\[
 \Cthreedd f=\Cthreedlog(1+uT).
\]
Evaluation at $T=0$ sends the specified relative unit to zero, so
$f(0)=0$.
The kernel of differentiation on $E\langle T\rangle$ consists of
constants: for each coefficient $c_n$ with $n\geq1$, the equation
$n c_n=0$ forces $c_n=0$.
The convergent restricted series
\[
 \log(1+uT)
   =\sum_{n\geq1}\frac{(-1)^{n+1}u^nT^n}{n}
\]
has the required derivative and constant term.
Thus $f=\log(1+uT)$.
For general $v\in1+uD$, substitution
$T=(v-1)/u\in D$ defines a continuous $W$-algebra map from
$W\langle T\rangle$ to $D$.
Naturality gives $\chi_D([v]_{\Cthreerel})=\log v$.

It remains to specify the coefficient in the product formula.
Over any commutative base ring, negative cyclic homology acts on
cyclic homology by ordinary shuffle plus the cyclic variable times
cyclic shuffle
\cite[\S2, Definition~2.1 and Remark~2.2]{BokstedtOttosen}.
The cyclic, negative cyclic, and periodic long exact sequences are
module sequences, with the usual sign on connecting homomorphisms
\cite[Proposition~2.3]{BokstedtOttosen}.
These statements precede, and do not require, the characteristic-two
hypotheses imposed in later sections of that reference.

Write the cyclic complex, with $h$ of degree $-2$, as
\[
 C^+=\bigl(E[h,h^{-1}]/hE[h]\bigr)\otimes_E C_\ast,
\]
where $C_\ast$ is the rationalized Hochschild chain complex.
The function $\log v$ is represented in column zero.
A negative cyclic representative of a unit has leading Hochschild
class $b^{-1}\otimes b$; its remaining terms contain positive powers
of $h$.
Multiplication on column zero kills all those terms, including the
cyclic-shuffle correction $h\,\mathrm{sh}'$.
The first multiplication remains in column zero, so the second
multiplication has the same property.
The degree-two shuffle has two terms.
After antisymmetrization both give the same wedge, canceling the
factor $2!$ in \eqref{c3:eq:finite-hkr}.
Changing a leading Hochschild representative by a boundary does not
change its resulting cyclic class.
The coefficient in \eqref{c3:eq:unit-character} is therefore one,
apart from the suspension sign already specified.
\end{proof}

\subsection{An ordinary residue criterion}\label{c3:sec:residue}

In this section, we detect differential classes on the formal quadric
of Setup~\ref{c3:setup:coefficient-ring} by a residue on the overlap of two
affine-plane charts.

\begin{lemma}
\label{c3:lem:ordinary-residue}
For $F\in E\langle Z\rangle$, the form initially defined where $X$ is
invertible by
\[
 \omega_F=\frac{\Cthreedd X}{X}\wedge\Cthreedd F
\]
extends to a global element of $\Omega_B^2$.
If $\omega_F$ belongs to the ordinary image $\Cthreedd\Omega_B^1$, then
$F(1)=F(-1)$.
\end{lemma}

\begin{proof}
Differentiating the equation of the quadric gives
$X\Cthreedd Y+Y\Cthreedd X+2Z\Cthreedd Z=0$.
On the locus where $Z$ is invertible,
\[
 \frac{\Cthreedd X}{X}\wedge\Cthreedd Z
   =-\frac{\Cthreedd X\wedge\Cthreedd Y}{2Z}.
\]
The loci where $X$ or $Z$ is invertible cover the quadric, since
$X=Z=0$ is incompatible with $XY+Z^2=1$.
This identity extends $\omega_F$ across its potential pole.

Before completion, let $U_+$ be the complement of the line
$X=0$, $Z=-1$, and let $U_-$ be the complement of the line
$X=0$, $Z=1$.
Put $a_+=1$ and $a_-=-1$.
Each $U_i$ is an affine plane with coordinates $X,t_i$, under
\begin{equation}\label{c3:eq:residue-charts}
 Z=a_i+Xt_i,\qquad
 Y=-t_i(2a_i+Xt_i).
\end{equation}
To see this, write $U_i=D(X)\cup D(Z+a_i)$.
The inverse coordinate is
$(Z-a_i)/X$ on $D(X)$ and $-Y/(Z+a_i)$ on $D(Z+a_i)$.
The equation $(Z-a_i)(Z+a_i)=-XY$ identifies these functions on the
overlap.
Their preimages cover the $(X,t_i)$-plane, since $2a_i$ is a unit.
Thus \eqref{c3:eq:residue-charts} is an isomorphism.
The intersection $U_+\cap U_-$ is $D(X)$.
After completion these charts are formal bidiscs, with overlap
$W\langle X,X^{-1},Z\rangle$.

On the $i$-th bidisc, define
\[
 q_i=-\bigl(F(Z)-F(a_i)\bigr)\frac{\Cthreedd X}{X}.
\]
Substitution from \eqref{c3:eq:residue-charts} makes
$F(Z)-F(a_i)$ divisible by $X$ in the restricted series algebra, so
$q_i$ is regular.
Differentiation gives $\Cthreedd q_i=\omega_F$.
On the overlap,
\begin{equation}\label{c3:eq:primitive-difference}
 q_+-q_-=\bigl(F(1)-F(-1)\bigr)\frac{\Cthreedd X}{X}.
\end{equation}

We use a coefficient calculation for closed one-forms on a bidisc.
Let
\[
 A(X,t)\Cthreedd X+D(X,t)\Cthreedd t
\]
be closed over $E\langle X,t\rangle$, and write
$A=\sum_{m,n\geq0}a_{m,n}X^mt^n$ and
$D=\sum_{m,n\geq0}b_{m,n}X^mt^n$.
The coefficients tend to zero.
For $c\in W$, substitution $t=c/X$ is defined on
$E\langle X,X^{-1}\rangle$.
The coefficient of $X^{-1}\Cthreedd X$ in the resulting form is the
convergent sum
\[
 \sum_{m\geq0}
       \bigl(a_{m,m+1}-b_{m+1,m}\bigr)c^{m+1}.
\]
Closedness means $\partial_t A=\partial_X D$, hence
\[
 (m+1)a_{m,m+1}=(m+1)b_{m+1,m}
 \qquad(m\geq0).
\]
Since $E$ has characteristic zero, each summand in the residue is
zero. This calculation does not divide an infinite series to
construct a primitive.

Assume now that $\omega_F=\Cthreedd q$ for a global
$q\in\Omega_B^1$.
The forms $q_i-q$ are closed on the two bidiscs.
Restrict them to the common slice $Z=0$.
In chart $i$, this restriction is the substitution
$t_i=-a_i/X$, with $-a_i\in W$.
Both restricted forms have residue zero by the coefficient
calculation.
Their difference, by \eqref{c3:eq:primitive-difference}, has residue
$F(1)-F(-1)$.
Therefore $F(1)=F(-1)$.
The argument concerns exactness in the ordinary quotient
$\Omega_B^2/\Cthreedd\Omega_B^1$.
\end{proof}

\subsection{Nonvanishing of the explicit class}\label{c3:sec:nonvanishing}

The goal of this section is to complete
Theorem~\ref{thm:explicit-three-dimensional} by detecting an actual
relative class on $B$ and passing to its absolute image.

\begin{proof}[Proof of Theorem~\ref{thm:explicit-three-dimensional}]
Proposition~\ref{c3:prop:integral} proves regularity and integral
vanishing.
It remains to prove that $c_{A_3}$ and $c_{\widehat A_3}$ have infinite
order.

The substitutions
\[
 x=uX,\qquad y=uY,\qquad z=uZ
\]
define a continuous homomorphism $\widehat A_3\to B$.
Indeed, every formal series in $x,y,z$ converges $u$-adically after
substitution, and the relation maps to
$3+u^2(XY+Z^2)=0$.
In $B$, put
\[
 a=\frac{1+uZ}{2},\qquad
 \lambda=1+uX,\qquad
 \mu=1+\frac{(1+uX)uY}{4},\qquad
 q=a(1-a)=\frac{1+3Z^2}{4}.
\]
These elements are units, and $a^2$ reduces to $1$ modulo $u$.
Consequently, the specified relative unit of
Proposition~\ref{c3:prop:relative-character} gives an actual class
\begin{equation}\label{c3:eq:relative-class}
 \gamma=[a^2]_{\Cthreerel}[\lambda][\mu]\in\CthreeK_3(B,uB).
\end{equation}
Its image in $\CthreeK_3(B)$ is
$[a^2][\lambda][\mu]=2[a][\lambda][\mu]=c_B$,
where $c_B$ is the image of either source class.

The map from relative to absolute $\CthreeK$-theory is injective in this
degree. To prove this using actual groups, put
\[
 C=B/uB=\CthreeF_3[X,Y,Z]/(XY+Z^2-1).
\]
A second pair of affine-plane charts describes $\CthreeSpec C$ as an
affine-line torsor over $\CthreePP^1_{\CthreeF_3}$.
The first chart has coordinates $X,t$ and equations
\[
 Z=1+Xt,\qquad Y=-t(2+Xt).
\]
It is the complement of the line $X=0$, $Z=-1$.
The second has coordinates $Y,s$ and equations
\[
 Z=-1-sY,\qquad X=-s(2+sY).
\]
It is the complement of the line $Y=0$, $Z=1$.
These charts are affine planes by the inverse-coordinate calculation
used in \eqref{c3:eq:residue-charts}.
Their omitted lines are disjoint, so the charts cover $\CthreeSpec C$.
On their intersection,
\begin{equation}\label{c3:eq:affine-line-transition}
 s=t^{-1},\qquad Y=-t^2X-2t.
\end{equation}
Thus projection to $t$ on the first chart and to $s$ on the second
glues to a morphism $\CthreeSpec C\to\CthreePP^1_{\CthreeF_3}$.
It is locally the affine-line projection, with affine-linear
transition of invertible slope $-t^2$.
In particular, it is a torsor under a line bundle.

Homotopy invariance on the two charts and their overlap, followed by
the two-open Mayer--Vietoris sequence, shows that pullback gives an
equivalence on $\CthreeK$-theory
\cite[\S7, Remark~3.5 and Proposition~4.1]{Quillen}.
The projective bundle formula
\cite[\S8, Theorem~2.1]{Quillen} identifies
$\CthreeK(\CthreePP^1_{\CthreeF_3})$ with $\CthreeK(\CthreeF_3)\oplus\CthreeK(\CthreeF_3)$.
For a finite field, the actual-group computation is
\[
 \CthreeK_{2i}(\CthreeF_q)=0,\qquad
 \CthreeK_{2i-1}(\CthreeF_q)=\CthreeZ/(q^i-1)
 \qquad(i\geq1);
\]
see \cite{QuillenFinite} and
\cite[IV, Theorem~1.12 and Corollary~1.13]{Weibel}.
Only its consequence $\CthreeK_4(\CthreeF_3)=0$ is needed here.
It follows that
\begin{equation}\label{c3:eq:actual-k-four}
 \CthreeK_4(C)=\CthreeK_4(\CthreeF_3)\oplus\CthreeK_4(\CthreeF_3)=0.
\end{equation}
The long exact sequence of the actual relative spectrum contains
\[
 \CthreeK_4(C)\longrightarrow\CthreeK_3(B,uB)\longrightarrow\CthreeK_3(B).
\]
Hence \eqref{c3:eq:actual-k-four} makes
$\CthreeK_3(B,uB)\to\CthreeK_3(B)$ injective.
This argument uses the actual fourth $\CthreeK$-group.

We detect $\gamma$ by its character.
Define the $3$-adic logarithm of $2$ by
\[
 \log2=\frac12\log4=\frac12\log(1+3),
 \qquad
 L(Z)=-\log2+\log(1+uZ).
\]
The logarithm series are additive on $1+uB$ by the formal power
series identity and their $u$-adic convergence.
Since $4\in1+uB$ and
$a^2=(1+uZ)^2/4$, this gives
$\log(a^2)=2L$.
This use of the logarithm remains valid at ramification index two.
Proposition~\ref{c3:prop:relative-character} identifies the top
component of $\chi_B(\gamma)$ with $\pm\tau$, where
\begin{equation}\label{c3:eq:character-form}
 \tau=2L\,\Cthreedlog\lambda\wedge\Cthreedlog\mu.
\end{equation}

For the differential calculation, put $t=uY/4$.
Then $q=1+uXt$, $\mu=1+\lambda t$, and
$\lambda q=1+uX\mu$.
Differentiating the last identity and wedging with $\Cthreedlog\mu$ gives,
where $X$ is invertible,
\begin{equation}\label{c3:eq:differential-identity}
 \Cthreedlog\lambda\wedge\Cthreedlog\mu
   =\frac{\Cthreedd X}{X}\wedge\Cthreedlog q
      -\Cthreedlog q\wedge\Cthreedlog\mu.
\end{equation}
More explicitly, the left-hand side plus
$\Cthreedlog q\wedge\Cthreedlog\mu$ is
\[
 \frac{u\,\Cthreedd X\wedge\Cthreedd\mu}{\lambda q}
   =\frac{u\,\Cthreedd X\wedge\Cthreedd t}{q}
   =\frac{\Cthreedd X}{X}\wedge\frac{\Cthreedd q}{q}.
\]
The form involving $\Cthreedd X/X$ is global by the extension calculation
in Lemma~\ref{c3:lem:ordinary-residue}, since $q$ depends only on $Z$.
Thus \eqref{c3:eq:differential-identity} is an identity of global regular
two-forms.

Choose $F\in E\langle Z\rangle$ satisfying
\begin{equation}\label{c3:eq:primitive-derivative}
 F'(Z)=\frac{6ZL(Z)}{1+3Z^2},
 \qquad\text{so that}\qquad
 \Cthreedd F=L\,\Cthreedlog q.
\end{equation}
Such an $F$ exists in the restricted series algebra.
With the absolute value normalized by $|3|=1/3$, both $L$ and
$(1+3Z^2)^{-1}$ converge for $|Z|<\sqrt3$.
Termwise integration divides each coefficient by a positive integer,
whose $3$-adic valuation grows at most logarithmically.
The integrated series therefore still has radius of convergence at
least $\sqrt3$, which is greater than one.
Equations \eqref{c3:eq:character-form} and
\eqref{c3:eq:differential-identity} now give the global identity
\begin{equation}\label{c3:eq:character-residue}
 \tau=
 2\frac{\Cthreedd X}{X}\wedge\Cthreedd F
       -2\Cthreedd\bigl(F\,\Cthreedlog\mu\bigr).
\end{equation}
By Lemma~\ref{c3:lem:ordinary-residue}, a nonzero difference
$F(1)-F(-1)$ makes the class of $\tau$ nonzero in the ordinary
quotient $\Omega_B^2/\Cthreedd\Omega_B^1$.

Only the odd part of $L$ contributes to this endpoint difference.
Since $u^2=-3$, that odd part is
\[
 u\sum_{k\geq0}\frac{(-3)^k Z^{2k+1}}{2k+1}.
\]
Multiply this series by $6Z/(1+3Z^2)$, integrate, and subtract the
values at $1$ and $-1$.
All series converge on a disc of radius greater than one, which
justifies termwise operations and regrouping by $j=k+l$.
The resulting formula is
\begin{equation}\label{c3:eq:endpoint-series}
 \frac{F(1)-F(-1)}{u}
   =12\sum_{j\geq0}\frac{(-3)^jS_j}{2j+3},
 \qquad
 S_j=\sum_{k=0}^j\frac1{2k+1}.
\end{equation}
Its $j=0$ term is $4$.

For completeness, an estimate for all indices controls the tail.
Normalize $v_3(3)=1$.
Every denominator in $S_j$ is at most $2j+1$, so
\[
 v_3(S_j)\geq-\lfloor\log_3(2j+1)\rfloor.
\]
For $j\geq1$, the valuation of the $j$-th summand in
\eqref{c3:eq:endpoint-series} is therefore at least
\[
 A_j=1+j-\lfloor\log_3(2j+1)\rfloor-v_3(2j+3).
\]
Direct substitution gives $A_1=1$, $A_2=2$, and $A_3=1$.
For $j\geq4$, we have
\begin{equation}\label{c3:eq:logarithmic-bound}
 \lfloor\log_3(2j+3)\rfloor\leq\lfloor j/2\rfloor.
\end{equation}
Equivalently, $2j+3<3^{\lfloor j/2\rfloor+1}$.
This holds for $j=4$ and $j=5$.
Advancing $j$ by two preserves the inequality because
$2j+7<3(2j+3)$.
Both $\lfloor\log_3(2j+1)\rfloor$ and $v_3(2j+3)$ are bounded by
$\lfloor\log_3(2j+3)\rfloor$.
Thus \eqref{c3:eq:logarithmic-bound} gives
\[
 A_j\geq1+j-2\lfloor j/2\rfloor\geq1.
\]
Moreover,
\[
 A_j\geq1+j-2\log_3(2j+3)\longrightarrow+\infty.
\]
The entire tail of \eqref{c3:eq:endpoint-series} consequently converges
in $3\CthreeZ_3$, and
\begin{equation}\label{c3:eq:endpoint-nonzero}
 \frac{F(1)-F(-1)}{u}\in4+3\CthreeZ_3.
\end{equation}
In particular, the endpoint difference is nonzero.

Equations \eqref{c3:eq:character-residue} and
\eqref{c3:eq:endpoint-nonzero}, together with
Lemma~\ref{c3:lem:ordinary-residue}, show that
$\chi_B(\gamma)$ has nonzero top component.
Its target is a $\CthreeQ_3$-vector space, so a torsion element cannot have
this image.
Hence $\gamma$ has infinite order in the actual group
$\CthreeK_3(B,uB)$.
The injection proved using \eqref{c3:eq:actual-k-four} shows that $c_B$
has infinite order in $\CthreeK_3(B)$.
As $c_B$ is the image of both $c_{\widehat A_3}$ and $c_{A_3}$, both
source classes have infinite order.

For either source ring $R$, Proposition~\ref{c3:prop:integral} gives
integral vanishing in $\CthreeK_3(R[1/x])$.
Since $R$ is a domain, $R[1/x]$ embeds in $\CthreeFrac(R)$, so the class
also vanishes in $\CthreeK_3(\CthreeFrac(R))$.
An element of an abelian group becomes zero after tensoring with
$\CthreeQ$ exactly when it is torsion.
Thus $c_R$ remains nonzero in $\CthreeK_3(R)\otimes_{\CthreeZ}\CthreeQ$, proving the
asserted failure of injectivity.
\end{proof}

\begin{remark}[Ordinary differential homology]
\label{c3:rem:ordinary-homology}
Every differential quotient used above is ordinary algebraic homology
after derived completion and rationalization.
No differential image is replaced by its closure.
The argument makes no assertion about the order of the class in an
individual finite Artin quotient, and uses no identification
\[
 \CthreeK_i(D)\otimes_{\CthreeZ}\CthreeQ_3
   =\pi_i\bigl(\CthreeK(D)^\wedge_3[1/3]\bigr).
\]
It also requires no injection from actual $\CthreeK$-theory into completed
$\CthreeK$-theory.

The distinction between an ordinary image and its closure is visible
on this quadric.
The global form
\[
 \omega_0=\frac{\Cthreedd X}{X}\wedge\Cthreedd Z
\]
is nonexact by Lemma~\ref{c3:lem:ordinary-residue}, applied to $F(Z)=Z$,
whose endpoint difference is $2$.
For every integer $j\geq0$, however, the one-form
\[
 \beta_j=-\bigl(Z-Z^{3^j}\bigr)\frac{\Cthreedd X}{X}
\]
is global and regular.
Indeed, $Z-Z^{3^j}$ vanishes at both $1$ and $-1$, and is divisible
by $Z^2-1=-XY$.
Its derivative is
\[
 \Cthreedd\beta_j
    =\omega_0-3^jZ^{3^j-1}\omega_0
    \longrightarrow\omega_0.
\]
Thus $\omega_0$ lies in the closure of the differential image while
remaining outside that image.
Passing to the Hausdorff quotient would discard the obstruction used
in the proof.
\end{remark}
